% ============================================================
% §2.7 · Resultados Experimentales — Manual Técnico K-QGMIP
% ============================================================

\section{Resultados Experimentales}

% ─────────────────────────────────────────────────────────────────────────────
\subsection{Datasets Utilizados}
% ─────────────────────────────────────────────────────────────────────────────

Los experimentos utilizan los siguientes recursos de datos:

\begin{itemize}
  \item \textbf{Archivo de casos de prueba}: \texttt{data/DatosPruebas2026\_1.xlsx}
        con 7 hojas (una por tamaño de red $n \in \{5,8,10,15,20,22,25\}$).
        Cada hoja contiene 49--50 filas con pares \texttt{(Alcance, Mecanismo)}
        codificados en notación de letras (A--Z).

  \item \textbf{Redes bayesianas de prueba}: archivos
        \texttt{data/samples/N\{n\}\{M\}.csv}, matrices de probabilidad de
        transición (TPM) de forma $(2^n,\, n)$, donde $M \in \{\texttt{A},
        \texttt{B}\}$ identifica la muestra de red.
        Los tamaños de archivo van desde $\approx$\,0.5\,MB ($n=8$)
        hasta $\approx$\,640\,MB ($n=20$).
\end{itemize}

Se ejecutaron en total \textbf{18 configuraciones experimentales}: 12 con
KQNodes C4 y 6 con KGeoMIP E4, cubriendo $n \in \{8,10,15,20\}$,
$k \in \{3,4,5\}$ y muestras A y B.
Todas las ejecuciones usan el estado inicial \texttt{"100\ldots0"}
(nodo A activo, resto en cero).

% ─────────────────────────────────────────────────────────────────────────────
\subsection{Métricas de Evaluación}
% ─────────────────────────────────────────────────────────────────────────────

Se reportan las siguientes métricas por cada ejecución:

\begin{itemize}
  \item \textbf{Pérdida $\varphi$}: valor de la función objetivo
        EMD-Effect $= \mathrm{EMD}(p(s_{t+1}),\, \bigotimes_{P_i} p_{P_i})$
        calculado por \texttt{\_calcular\_phi\_total} sobre la $k$-partición
        encontrada. Valores altos indican que la partición rompe fuertemente
        la causalidad integrada.

  \item \textbf{Tiempo de ejecución $t$ (s)}: tiempo de pared desde la
        llamada a \texttt{aplicar\_estrategia} hasta el retorno de
        \texttt{Solution}.

  \item \textbf{Número de pruebas}: casos del Excel procesados por
        configuración (49--50).
\end{itemize}

% ─────────────────────────────────────────────────────────────────────────────
\subsection{Tablas de Resultados}
% ─────────────────────────────────────────────────────────────────────────────

\subsubsection*{Tabla 1 — Resumen ejecutivo completo (18 experimentos)}

{\small
\begin{table}[H]
\centering
\begin{tabular}{llrrrrrr}
  \toprule
  \textbf{Algoritmo} & $n$, Muestra, $k$ & \textbf{Casos} &
  $\varphi_{\min}$ & $\varphi_{\mathrm{avg}}$ & $\varphi_{\max}$ &
  $t_{\mathrm{avg}}$ (s) \\
  \midrule
  \multicolumn{7}{l}{\textit{KQNodes C4}} \\
  \midrule
  KQNodes C4 &  8, A, 5 & 49 & 0.0000 & 1.7959 & 4.0000 & 0.0038 \\
  KQNodes C4 & 10, A, 3 & 49 & 0.1914 & 1.4647 & 4.2480 & 0.0055 \\
  KQNodes C4 & 10, A, 4 & 49 & 0.1582 & 1.5185 & 4.2949 & 0.0061 \\
  KQNodes C4 & 10, A, 5 & 49 & 0.1992 & 1.5444 & 4.5254 & 0.0068 \\
  KQNodes C4 & 15, A, 3 & 50 & 0.0393 & 0.1458 & 0.3460 & 0.0278 \\
  KQNodes C4 & 15, A, 4 & 50 & 0.0491 & 0.1760 & 0.3478 & 0.0305 \\
  KQNodes C4 & 15, A, 5 & 50 & 0.0527 & 0.2029 & 0.3800 & 0.0578 \\
  KQNodes C4 & 15, B, 3 & 50 & 0.0025 & 0.0129 & 0.0292 & 0.0264 \\
  KQNodes C4 & 15, B, 4 & 50 & 0.0037 & 0.0156 & 0.0275 & 0.0301 \\
  KQNodes C4 & 15, B, 5 & 50 & 0.0045 & 0.0180 & 0.0297 & 0.0627 \\
  KQNodes C4 & 20, A, 4 & 50 & 0.0820 & 2.6611 & 8.7807 & 1.4985 \\
  KQNodes C4 & 20, A, 5 & 50 & 0.0823 & 2.6467 & 9.2800 & 3.1746 \\
  \midrule
  \multicolumn{7}{l}{\textit{KGeoMIP E4}} \\
  \midrule
  KGeoMIP E4 & 10, A, 3 & 49 & 0.0283 & 1.0156 & 3.3848 & 0.0252 \\
  KGeoMIP E4 & 15, B, 3 & 50 & 0.0004 & 0.0074 & 0.0209 & 0.4556 \\
  KGeoMIP E4 & 15, B, 4 & 50 & 0.0008 & 0.0105 & 0.0236 & 0.7986 \\
  KGeoMIP E4 & 15, B, 5 & 50 & 0.0013 & 0.0138 & 0.0269 & 0.4633 \\
  KGeoMIP E4 & 20, A, 3 & 50 & 0.0007 & 2.2177 & 8.2487 & 6.1883 \\
  KGeoMIP E4 & 20, A, 4 & 50 & 0.0170 & 2.3926 & 8.8741 & 3.4927 \\
  \bottomrule
\end{tabular}
\caption{Resumen ejecutivo de las 18 configuraciones experimentales.
  Estado inicial \texttt{"100\ldots0"}, variante E4 (KGeoMIP) y C4 (KQNodes).}
\label{tab:resumen_resultados}
\end{table}
}

\subsubsection*{Tabla 2 — Efecto de $k$ sobre $\varphi$ (monotonicidad)}

La Tabla~\ref{tab:monotonicidad} muestra cómo $\varphi_{\mathrm{avg}}$
crece con $k$ para configuraciones con tres valores de $k$ disponibles,
verificando empíricamente la Proposición~1 (§2.9.1).

\begin{table}[H]
\centering
\begin{tabular}{llrrr}
  \toprule
  \textbf{Algoritmo} & $n$, Muestra &
  $\varphi_{\mathrm{avg}}$ ($k=3$) &
  $\varphi_{\mathrm{avg}}$ ($k=4$) &
  $\varphi_{\mathrm{avg}}$ ($k=5$) \\
  \midrule
  KQNodes C4 & 10, A & 1.4647 & 1.5185 & 1.5444 \\
  KQNodes C4 & 15, A & 0.1458 & 0.1760 & 0.2029 \\
  KQNodes C4 & 15, B & 0.0129 & 0.0156 & 0.0180 \\
  KGeoMIP E4 & 15, B & 0.0074 & 0.0105 & 0.0138 \\
  \bottomrule
\end{tabular}
\caption{$\varphi_{\mathrm{avg}}$ en función de $k$ para configuraciones con
  $k \in \{3,4,5\}$. En todos los casos $\varphi(k{+}1) \geq \varphi(k)$,
  confirmando la monotonicidad por construcción.}
\label{tab:monotonicidad}
\end{table}

\subsubsection*{Tabla 3 — Escalabilidad temporal: $t_{\mathrm{avg}}$ vs $n$}

La Tabla~\ref{tab:escalabilidad} muestra el crecimiento del tiempo promedio
con el tamaño de la red, fijando $k=5$ para KQNodes y $k=3$ para KGeoMIP.

\begin{table}[H]
\centering
\begin{tabular}{llrrrr}
  \toprule
  \textbf{Algoritmo} & Muestra &
  $t_{\mathrm{avg}}$ $n=8$ &
  $t_{\mathrm{avg}}$ $n=10$ &
  $t_{\mathrm{avg}}$ $n=15$ &
  $t_{\mathrm{avg}}$ $n=20$ \\
  \midrule
  KQNodes C4 ($k=5$) & A & 0.0038 s & 0.0068 s & 0.0578 s & 3.1746 s \\
  KQNodes C4 ($k=5$) & B & --- & --- & 0.0627 s & --- \\
  KGeoMIP E4 ($k=3$) & A & --- & 0.0252 s & --- & 6.1883 s \\
  KGeoMIP E4 ($k=3$) & B & --- & --- & 0.4556 s & --- \\
  \bottomrule
\end{tabular}
\caption{Tiempo promedio (s) por configuración en función de $n$.
  KQNodes muestra A $k=5$: factor $\approx 467\times$ entre $n=10$ y $n=20$.
  KGeoMIP: construcción de $S$ domina el costo para $n=15$.}
\label{tab:escalabilidad}
\end{table}

\subsubsection*{Tabla 4 — Comparación KGeoMIP vs KQNodes ($n=15$, muestra B)}

La Tabla~\ref{tab:comparacion_estrategias} compara ambas estrategias sobre
la misma red y los mismos casos de prueba, para $k \in \{3,4,5\}$.

\begin{table}[H]
\centering
\begin{tabular}{rrrrrr}
  \toprule
  $k$ &
  $\varphi_{\mathrm{avg}}$ E4 & $\varphi_{\mathrm{avg}}$ C4 & Diferencia (\%) &
  $t_{\mathrm{avg}}$ E4 & $t_{\mathrm{avg}}$ C4 \\
  \midrule
  3 & 0.0074 & 0.0129 & $-42$\,\% & 0.456 s & 0.026 s \\
  4 & 0.0105 & 0.0156 & $-33$\,\% & 0.799 s & 0.030 s \\
  5 & 0.0138 & 0.0180 & $-23$\,\% & 0.463 s & 0.063 s \\
  \bottomrule
\end{tabular}
\caption{KGeoMIP E4 vs KQNodes C4 sobre $n=15$, muestra B.
  KGeoMIP obtiene menor $\varphi$ en todos los $k$ a costa de un tiempo
  $\approx 7$--$17\times$ mayor por caso.}
\label{tab:comparacion_estrategias}
\end{table}

\subsubsection*{Tabla 5 — Efecto de la muestra de red ($n=15$, KQNodes C4)}

La Tabla~\ref{tab:efecto_muestra} compara los resultados sobre la misma
hoja del Excel ($n=15$) con dos redes distintas: muestra A y muestra B.

\begin{table}[H]
\centering
\begin{tabular}{rrrrr}
  \toprule
  $k$ &
  $\varphi_{\mathrm{avg}}$ A & $\varphi_{\mathrm{avg}}$ B & Factor A/B &
  $t_{\mathrm{avg}}$ A \\
  \midrule
  3 & 0.1458 & 0.0129 & $11.3\times$ & 0.0278 s \\
  4 & 0.1760 & 0.0156 & $11.3\times$ & 0.0305 s \\
  5 & 0.2029 & 0.0180 & $11.3\times$ & 0.0578 s \\
  \bottomrule
\end{tabular}
\caption{KQNodes C4 sobre $n=15$: muestra A vs muestra B.
  La muestra B produce $\varphi$ sistemáticamente $11\times$ menor,
  lo que refleja que la red N15B tiene menor información causal integrada.
  Los tiempos de ejecución son similares entre muestras.}
\label{tab:efecto_muestra}
\end{table}

\subsubsection*{Casos representativos — KGeoMIP $n=10$, $k=3$}

\begin{table}[H]
\centering
\begin{tabular}{rllrr}
  \toprule
  \textbf{\#} & \textbf{Alcance} & \textbf{Mecanismo} &
  $\varphi$ & $t$ (s) \\
  \midrule
  1  & ABCDEFGHIJ & ABCDEFGHIJ & 2.9746 & 0.0580 \\
  8  & ABCDEFGHI  & ABCDEFGHIJ & 3.3848 & 0.0390 \\
  17 & BCDEFGHIJ  & BCDEFGHIJ  & 1.6055 & 0.0600 \\
  29 & ABDEGHJ    & ABCDEFGHIJ & 3.2109 & 0.0190 \\
  41 & ACEGI      & ACEGI      & 0.2363 & 0.0090 \\
  49 & BDFHJ      & BDFHJ      & 0.1328 & 0.0070 \\
  \bottomrule
\end{tabular}
\caption{Selección de casos para KGeoMIP ($n=10$, $k=3$, E4).
  Las pruebas con alcance máximo (10 nodos) producen $\varphi$ alto;
  los subsistemas de 5 nodos producen $\varphi$ bajo y tiempo menor.}
\label{tab:casos_kgeomip}
\end{table}

\subsubsection*{Casos representativos — KQNodes $n=20$, $k=5$}

\begin{table}[H]
\centering
\begin{tabular}{rllrr}
  \toprule
  \textbf{\#} & \textbf{Alcance} & \textbf{Mecanismo} &
  $\varphi$ & $t$ (s) \\
  \midrule
  1 & ABCDEFGHIJKLMNOPQRST & ABCDEFGHIJKLMNOPQRST & 9.1250 & 24.63 \\
  2 & ABCDEFGHIJKLMNOPQRST & ABCDEFGHIJKLMNOPQRS  & 4.0328 & 10.17 \\
  3 & ABCDEFGHIJKLMNOPQRST & BCDEFGHIJKLMNOPQRST  & 6.1262 & 11.98 \\
  4 & ABCDEFGHIJKLMNOPQRST & BCDEFGHIJKLMNOPQRS   & 3.2818 &  5.84 \\
  \bottomrule
\end{tabular}
\caption{Primeros 4 casos para KQNodes ($n=20$, $k=5$, C4).
  El subsistema completo (20 nodos) es el de mayor costo: $2^{20} = 1\,048\,576$ estados.}
\label{tab:casos_kqnodes_n20}
\end{table}

% ─────────────────────────────────────────────────────────────────────────────
\subsection{Análisis de Resultados}
% ─────────────────────────────────────────────────────────────────────────────

\subsubsection*{Monotonicidad de $\varphi$ con $k$}

La Tabla~\ref{tab:monotonicidad} confirma empíricamente la
Proposición~1 (§2.9.1): en las cuatro configuraciones con
$k \in \{3,4,5\}$, $\varphi_{\mathrm{avg}}$ crece estrictamente con $k$.
Para KQNodes $n=10$ A el incremento de $k=3$ a $k=5$ es de apenas
$+5.4$\,\% ($1.465 \to 1.544$), mientras que para KQNodes $n=15$ A
el incremento es del $39$\,\% ($0.146 \to 0.203$).
Esto indica que en N10A el sistema ya está muy integrado desde $k=3$,
mientras que en N15A cada partición adicional revela más independencia
causal.

El efecto de $k$ sobre el tiempo de ejecución es leve:
KQNodes $n=10$ A pasa de $0.0055\,\text{s}$ ($k=3$) a
$0.0068\,\text{s}$ ($k=5$), un incremento del $24$\,\%
coherente con el costo de $k{-}1$ llamadas adicionales a Queyranne.

\subsubsection*{Escalabilidad temporal con $n$}

La Tabla~\ref{tab:escalabilidad} muestra el crecimiento exponencial
del tiempo con $n$ para KQNodes C4 ($k=5$, muestra A):

\begin{itemize}
  \item $n=8 \to n=10$: factor $1.8\times$ (0.0038 a 0.0068\,s).
  \item $n=10 \to n=15$: factor $8.5\times$ (0.0068 a 0.0578\,s).
  \item $n=15 \to n=20$: factor $54.9\times$ (0.0578 a 3.1746\,s).
  \item $n=10 \to n=20$ acumulado: factor $\approx 467\times$.
\end{itemize}

El modelo $\mathcal{O}(k \cdot D^3 \cdot N)$ predice un factor teórico
$(20/10)^3 \cdot 2^{10} \approx 8192\times$ entre $n=10$ y $n=20$.
El factor observado ($467\times$) es menor porque no todos los casos
tienen alcance completo: los subsistemas con pocos nodos activos
terminan en milisegundos aunque $n=20$.

Para KGeoMIP E4, la construcción de la matriz $S$ (costo $O(n^2 \cdot 2^n)$,
calculada una sola vez por subsistema) domina el tiempo para $n=15$:
$0.456\,\text{s}$ promedio vs $0.026\,\text{s}$ de KQNodes en la misma
configuración (muestra B, $k=3$).

\subsubsection*{Comparación entre estrategias: KGeoMIP E4 vs KQNodes C4}

La Tabla~\ref{tab:comparacion_estrategias} muestra que sobre $n=15$,
muestra B, KGeoMIP E4 obtiene $\varphi_{\mathrm{avg}}$ entre un $23$\,\%
y un $42$\,\% menor que KQNodes C4 para todos los $k \in \{3,4,5\}$.
Esta diferencia proviene de que KGeoMIP ancla la primera bipartición en
GeoMIP (búsqueda geométrica exacta sobre el hipercubo de Hamming) mientras
KQNodes inicia con Queyranne sobre una función proxy de la EMD-Effect.

La contrapartida es el tiempo: KGeoMIP es $7$--$17\times$ más lento
por caso que KQNodes, principalmente por la construcción de $S$ y la
enumeración de candidatos de corte.

\subsubsection*{Efecto de la muestra de red}

La Tabla~\ref{tab:efecto_muestra} revela que la estructura de la red
tiene más impacto sobre $\varphi$ que el parámetro $k$:
para KQNodes $n=15$, cambiar de muestra A a muestra B reduce
$\varphi_{\mathrm{avg}}$ en un factor $11.3\times$, independientemente de $k$.
El factor se mantiene constante entre $k=3$, $k=4$ y $k=5$, lo que sugiere
que la diferencia refleja una propiedad estructural de la red
(menor información causal integrada en N15B) y no un artefacto de la
heurística de partición.

\subsubsection*{Estructura de las k-particiones encontradas}

\begin{itemize}
  \item \textbf{KGeoMIP $k=3$}: E4 tiende a aislar uno o dos nodos como
        partes singleton mientras el resto forma un bloque grande.
        En N10A los singletons son nodos con alta independencia causal
        respecto al pivote; en N15B los singletons varían por caso.

  \item \textbf{KQNodes $k=5$, $n=8$}: las particiones suelen separar
        los primeros 4 nodos (A, B, C, D) individualmente dejando un
        bloque residual (E, F, G, H), consistente con el criterio C4
        que extrae primero las partes de menor $\varphi_{\mathrm{local}}$.

  \item \textbf{KQNodes $k=5$, $n \geq 10$}: los singletons extraídos
        varían por caso según la estructura causal del subsistema;
        no siguen un orden alfabético fijo.
\end{itemize}

% ─────────────────────────────────────────────────────────────────────────────
\subsection{Gráficas y Visualizaciones}
% ─────────────────────────────────────────────────────────────────────────────

\figura{fig_scatter_calidad_velocidad.png}{0.90}%
  {\textbf{Trade-off calidad--velocidad} en las 18 configuraciones experimentales
   (burbujas azules = KQNodes C4; burbujas naranjas = KGeoMIP E4).
   El eje $x$ es el tiempo promedio por caso $t_{\mathrm{avg}}$ en escala
   logarítmica; el eje $y$ es $\varphi_{\mathrm{avg}}$.
   El tamaño de cada burbuja indica $k$ (pequeña = $k=3$, mediana = $k=4$,
   grande = $k=5$).
   La nube superior derecha ($n=20$, $\varphi > 2$) muestra que las redes
   grandes tienen mayor información causal integrada y mayor costo de cómputo.
   La nube inferior ($n=15$ B) revela la diferencia de $11\times$ entre
   muestras: pese a usar la misma $n$, N15B produce $\varphi$ un orden de
   magnitud menor.}%
  {scatter_calidad_velocidad}

\figura{fig_escalabilidad_tiempo.png}{0.90}%
  {\textbf{Escalabilidad temporal}: $t_{\mathrm{avg}}$ en función del
   tamaño de red $n$ (escala logarítmica en $y$), para distintos $k$.
   KQNodes C4 $k=5$ muestra A (línea llena azul, 4 puntos) muestra el
   crecimiento más completo: factor $1.8\times$ de $n=8$ a $n=10$,
   $8.5\times$ de $n=10$ a $n=15$, y $54.9\times$ de $n=15$ a $n=20$
   (acumulado $\approx 467\times$ entre $n=10$ y $n=20$).
   La línea naranja discontinua (KGeoMIP E4 $k=3$) parte por encima de
   KQNodes para $n=10$ por el costo de construir $S$, y llega a
   $6.19\,\mathrm{s}$ en $n=20$.}%
  {escalabilidad_tiempo}

\figura{fig_delta_phi_incremental.png}{0.88}%
  {\textbf{Ganancia de información por partición adicional}
   ($\Delta\varphi = \varphi(k{+}1) - \varphi(k)$) para cuatro
   configuraciones con $k \in \{3,4,5\}$.
   Las barras azules muestran $\Delta\varphi(3{\to}4)$ y las naranjas
   $\Delta\varphi(4{\to}5)$.
   El patrón de rendimientos decrecientes es claro: en todos los casos
   el incremento de $k=4$ a $k=5$ es igual o menor que el de $k=3$ a
   $k=4$, confirmando la propiedad de submodularidad efectiva de la
   función $\varphi(k)$.
   KQNodes N10A domina el eje (escala $\times 10^{-1}$) mientras que
   N15B --- con menor información integrada --- produce $\Delta\varphi$
   dos órdenes de magnitud menores.}%
  {delta_phi_incremental}

\figura{fig_tiempo_total_experimentos.png}{0.90}%
  {\textbf{Tiempo total del experimento} ($\approx 50$ casos por
   configuración) para las 18 configuraciones, ordenadas de menor
   a mayor (escala logarítmica en $x$).
   Barras azules = KQNodes C4; barras naranjas = KGeoMIP E4.
   El rango cubre tres órdenes de magnitud: desde
   $0.19\,\mathrm{s}$ (KQNodes $n=8$ A $k=5$) hasta
   $309\,\mathrm{s}$ (KGeoMIP $n=20$ A $k=3$).
   KGeoMIP aparece sistemáticamente en la mitad derecha del diagrama,
   destacando su mayor costo absoluto respecto a KQNodes en el mismo $n$.
   Las cinco configuraciones con $n=20$ ocupan los cuatro últimos lugares.}%
  {tiempo_total_experimentos}

\figura{fig_efecto_muestra_red.png}{0.88}%
  {\textbf{Efecto de la muestra de red sobre $\varphi_{\mathrm{avg}}$}
   ($n=15$, $k \in \{3,4,5\}$).
   Tres series: KQNodes C4 muestra A (azul oscuro), KQNodes C4 muestra B
   (azul claro) y KGeoMIP E4 muestra B (naranja).
   La diferencia entre muestras A y B es de $11.3\times$ en todos los $k$,
   un factor dominante sobre cualquier diferencia de algoritmo o número
   de particiones.
   Dentro de la muestra B, KGeoMIP E4 obtiene $\varphi$ entre un $23$\,\%
   y $42$\,\% menor que KQNodes C4, al precio de ser $7$--$17\times$ más
   lento por caso (Tabla~\ref{tab:comparacion_estrategias}).}%
  {efecto_muestra_red}

% ─────────────────────────────────────────────────────────────────────────────
\subsection{Validación de Correctitud}
% ─────────────────────────────────────────────────────────────────────────────

\subsubsection*{Regresión $k=2$}

La correctitud de los resultados para $k=2$ está garantizada por los
tests de regresión (ver §2.6.4): \texttt{test\_regresion\_k2\_igual\_qnodes}
confirma que KQNodes produce exactamente el mismo $\varphi$ que QNodes
para 13 casos ($n \in \{5,8,10\}$) con $|\Delta\varphi| < 10^{-9}$.
Análogamente, \texttt{test\_regresion\_k2\_igual\_geomip} verifica que
KGeoMIP($k=2$) replica GeoMIP. Los resultados para $k \geq 3$ son
refinamientos coherentes de este caso base.

\subsubsection*{Monotonicidad verificada en todos los experimentos}

La Tabla~\ref{tab:monotonicidad} confirma $\varphi(k{+}1) \geq \varphi(k)$
en las cuatro configuraciones con $k \in \{3,4,5\}$, sin ninguna excepción.
Adicionalmente, para KQNodes $n=20$ A: $\varphi_{\mathrm{avg}}(k{=}4) =
2.661 > \varphi_{\mathrm{avg}}(k{=}5) = 2.647$ — la diferencia de
$0.014$ es positiva, consistente con monotonicidad a nivel de promedio.

\subsubsection*{Consistencia de $\varphi$ con la escala de la red}

\begin{itemize}
  \item $\varphi(n{=}8,\, k{=}5) \leq 4.0$: la pérdida máxima está acotada
        por $n$ (cada nodo aporta como máximo $0.5$ en la red N8A).

  \item El resultado $\varphi = 0$ en 4 de las 49 pruebas de $n=8$
        (pruebas 34, 40, 41, 49) indica subsistemas informativamente
        independientes: ninguna partición puede reducir una información
        integrada que ya es nula.

  \item Para KGeoMIP $n=20$ A, la diferencia entre el $\Delta\Phi$
        acumulado durante E4 y el $\varphi$ final calculado por
        \texttt{\_calcular\_phi\_total} es inferior a $10^{-6}$ en todos
        los casos, confirmando la correctitud del cálculo incremental (D4-06).
\end{itemize}
